% Auto-split to keep each TeX source < 600 lines.

\paragraph{Solvable information geometry of bin-fold: escort temperature family, divergence spectra, and two-scale corrections for gauge density}
The two-state uniform asymptotics of Theorem \ref{thm:fold-bin-two-state-asymptotic} shows that at exponential leading order, fiber multiplicity depends only on the last bit \(w_m\).
This degeneracy compresses several commonly used information-geometric functionals (escort temperature families, Rényi/Csiszár divergence families, and gauge volume density) uniformly into the same set of explicit constants, with uniform exponential convergence rates.

\begin{proposition}[Last-bit sufficiency of bin-fold pushforward distribution: total variation collapse and exponential tilting model]\label{prop:fold-bin-lastbit-sufficient-tv}
Adopting the notation of \(\mu_m\) from Corollary \ref{cor:fold-bin-two-point-limit-law}, and letting \(\ell(w):=w_m\in\{0,1\}\) denote the last-bit indicator,
partition \(X_m\) by last bit into two blocks
\[
A_i:=\{w\in X_m:\ \ell(w)=i\}\qquad(i\in\{0,1\}),
\]
and define the last-bit conditioned approximation (block-wise uniformization)
\[
\overline\mu_m(w):=\frac{\mu_m(A_{\ell(w)})}{|A_{\ell(w)}|},\qquad w\in X_m.
\]
Then there exists an absolute constant \(C>0\) such that for all sufficiently large \(m\),
\[
D_{\mathrm{TV}}(\mu_m,\overline\mu_m)
:=\frac12\sum_{w\in X_m}\bigl|\mu_m(w)-\overline\mu_m(w)\bigr|
\ \le\ C\Bigl(\frac{\varphi}{2}\Bigr)^m.
\]
Moreover, there exists \(\beta_m\in\RR\) such that \(\overline\mu_m\) can be written as an exponential tilt of the stable-layer uniform distribution \(u_m(w)=1/|X_m|\):
\[
\overline\mu_m(w)=\frac{1}{Z_m}\,u_m(w)\,e^{-\beta_m\ell(w)},
\qquad
Z_m:=\sum_{w\in X_m}u_m(w)e^{-\beta_m\ell(w)},
\]
and \(\beta_m\to\log\varphi\).
\end{proposition}
\begin{proof}
By the uniform bound \eqref{eq:fold-bin-two-state-asymptotic} of Theorem \ref{thm:fold-bin-two-state-asymptotic}, within each last-bit block \(A_i\),
\(
d^{\mathrm{bin}}_m(w)
\)
differs from its block-wise average by \(O(1)\) (uniformly in \(w\in A_i\)), thus
\[
\sup_{w\in A_i}\bigl|\mu_m(w)-\overline\mu_m(w)\bigr|
=\sup_{w\in A_i}\Bigl|\frac{d^{\mathrm{bin}}_m(w)}{2^m}-\frac{1}{|A_i|}\sum_{w'\in A_i}\frac{d^{\mathrm{bin}}_m(w')}{2^m}\Bigr|
\ \le\ \frac{C_1}{2^m}.
\]
Summing over both blocks yields
\[
\sum_{w\in X_m}\bigl|\mu_m(w)-\overline\mu_m(w)\bigr|
\le
|X_m|\cdot \frac{C_1}{2^m}
=O\!\Bigl(\Bigl(\frac{\varphi}{2}\Bigr)^m\Bigr),
\]
where we used \(|X_m|=F_{m+2}=\Theta(\varphi^m)\), thus obtaining the shown total variation bound.
\par
For the exponential tilt representation, note that \(\overline\mu_m\) is constant on each \(A_i\): \(\overline\mu_m|_{A_i}=\mu_m(A_i)/|A_i|\).
Let
\[
e^{-\beta_m}:=\frac{\mu_m(A_1)/|A_1|}{\mu_m(A_0)/|A_0|},
\]
then for any \(w\in A_i\) we have \(\overline\mu_m(w)\propto e^{-\beta_m i}\), so there exists a normalization constant \(Z_m\) such that
\(
\overline\mu_m(w)=Z_m^{-1}u_m(w)e^{-\beta_m\ell(w)}
\).
\par
Finally, by the block mass limits of Corollary \ref{cor:fold-bin-two-point-limit-law}
\(\mu_m(A_0)\to \varphi/\sqrt5\), \(\mu_m(A_1)\to 1/(\varphi\sqrt5)\), and \(|A_0|=F_{m+1}\), \(|A_1|=F_m\) giving
\(|A_1|/|A_0|\to \varphi^{-1}\), we obtain
\[
e^{-\beta_m}
=\frac{\mu_m(A_1)}{\mu_m(A_0)}\cdot\frac{|A_0|}{|A_1|}
\longrightarrow
\frac{\varphi^{-2}}{1}\cdot \varphi
=\varphi^{-1},
\]
therefore \(\beta_m\to\log\varphi\). This completes the proof.
\end{proof}

\begin{proposition}[Last-bit limit law of escort temperature family and fiber information constant (bin-fold)]\label{prop:fold-bin-escort-lastbit}
For any \(q\ge 0\), define the escort distribution
\[
\pi_{m,q}(w):=\frac{\bigl(d^{\mathrm{bin}}_m(w)\bigr)^q}{S_q^{\mathrm{bin}}(m)},\qquad w\in X_m,
\]
and let the random variable \(W^{(q)}_m\sim\pi_{m,q}\).
Set
\[
\theta(q):=\frac{1}{1+\varphi^{q+1}}.
\]
Then as \(m\to\infty\),
\[
\mathbb{P}\bigl(W^{(q)}_m(m)=1\bigr)\to \theta(q),
\qquad
\mathbb{P}\bigl(W^{(q)}_m(m)=0\bigr)\to 1-\theta(q),
\]
and for fixed \(q\) the convergence rate is \(O\bigl((\varphi/2)^m\bigr)\).

Furthermore, define the \(q\)-escort fiber information
\[
\kappa^{\mathrm{bin}}_m(q):=\mathbb{E}_{\pi_{m,q}}\bigl[\log d^{\mathrm{bin}}_m(W^{(q)}_m)\bigr],
\]
then there exists a uniform expansion
\[
\kappa^{\mathrm{bin}}_m(q)=m\log\!\Bigl(\frac{2}{\varphi}\Bigr)-\theta(q)\log\varphi+O\bigl((\varphi/2)^m\bigr)
=m\log\!\Bigl(\frac{2}{\varphi}\Bigr)-\frac{\log\varphi}{1+\varphi^{q+1}}+O\bigl((\varphi/2)^m\bigr),
\]
where the error is uniform in \(m\) for fixed \(q\).
\end{proposition}

\begin{proof}
By definition,
\[
\mathbb{P}\bigl(W^{(q)}_m(m)=1\bigr)
=\sum_{w_m=1}\pi_{m,q}(w)
=\frac{\sum_{w_m=1}\bigl(d^{\mathrm{bin}}_m(w)\bigr)^q}{S_q^{\mathrm{bin}}(m)}.
\]
By the uniform leading term of Theorem \ref{thm:fold-bin-two-state-asymptotic}, we obtain
\(
\bigl(d^{\mathrm{bin}}_m(w)\bigr)^q=\bigl(\frac{2}{\varphi}\bigr)^{qm}\varphi^{-qw_m}\bigl(1+O(q(\varphi/2)^m)\bigr)
\)
uniformly on \(w\in X_m\).
Therefore
\[
\sum_{w_m=1}\bigl(d^{\mathrm{bin}}_m(w)\bigr)^q
=F_m\Bigl(\frac{2}{\varphi}\Bigr)^{qm}\varphi^{-q}\Bigl(1+O\bigl(q(\varphi/2)^m\bigr)\Bigr),
\]
while Proposition \ref{prop:fold-bin-power-sum-asymptotic} gives
\(
S_q^{\mathrm{bin}}(m)=\bigl(\frac{2}{\varphi}\bigr)^{qm}\bigl(F_{m+1}+\varphi^{-q}F_m\bigr)\bigl(1+O(q(\varphi/2)^m)\bigr)
\).
Dividing the two yields
\[
\mathbb{P}\bigl(W^{(q)}_m(m)=1\bigr)
=\frac{F_m\varphi^{-q}}{F_{m+1}+\varphi^{-q}F_m}+O\bigl((\varphi/2)^m\bigr)
\longrightarrow\ \frac{1}{1+\varphi^{q+1}},
\]
where we used \(F_{m+1}/F_m\to\varphi\).
The other endpoint probability follows by complement.

For \(\kappa^{\mathrm{bin}}_m(q)\), by the uniform logarithmic expansion of Theorem \ref{thm:fold-bin-two-state-asymptotic}
\[
\log d^{\mathrm{bin}}_m(w)
=m\log\!\Bigl(\frac{2}{\varphi}\Bigr)-w_m\log\varphi+O\bigl((\varphi/2)^m\bigr),
\qquad (w\in X_m),
\]
taking expectation under \(\pi_{m,q}\) and substituting \(\mathbb{P}(W^{(q)}_m(m)=1)=\theta(q)+O((\varphi/2)^m)\) yields the conclusion.
\end{proof}

\begin{proposition}[Escort-escort KL limit between temperatures (bin-fold)]\label{prop:fold-bin-escort-escort-kl}
For \(p,q\ge 0\), let \(\pi_{m,q}\) be as in Proposition \ref{prop:fold-bin-escort-lastbit}.
Then the limit
\[
\lim_{m\to\infty}D_{\mathrm{KL}}(\pi_{m,q}\Vert \pi_{m,p})
\]
exists and can be written as a closed form of one-dimensional Bernoulli KL. Setting \(\theta(\cdot)\) as defined in Proposition \ref{prop:fold-bin-escort-lastbit}, we have
\[
\lim_{m\to\infty}D_{\mathrm{KL}}(\pi_{m,q}\Vert \pi_{m,p})
=\theta(q)\log\Bigl(\frac{\theta(q)}{\theta(p)}\Bigr)+\bigl(1-\theta(q)\bigr)\log\Bigl(\frac{1-\theta(q)}{1-\theta(p)}\Bigr).
\]
\end{proposition}

\begin{proof}
By Theorem \ref{thm:fold-bin-two-state-asymptotic}, for any fixed \(q\ge 0\),
\[
\pi_{m,q}(w)=\frac{\varphi^{-qw_m}}{F_{m+1}+\varphi^{-q}F_m}\Bigl(1+O\bigl((\varphi/2)^m\bigr)\Bigr),
\qquad (w\in X_m),
\]
where the error is uniform in \(w\); similarly for \(p\).
Therefore, conditioned on \(w_m\), both \(\pi_{m,q}\) and \(\pi_{m,p}\) are asymptotically uniform within that sector, and the ratio \(\pi_{m,q}(w)/\pi_{m,p}(w)\) is also asymptotically constant within the sector.
Decomposing KL divergence by the two sectors \(w_m\in\{0,1\}\) and using
\(\mathbb{P}_{\pi_{m,q}}(w_m=1)\to \theta(q)\), \(\mathbb{P}_{\pi_{m,p}}(w_m=1)\to \theta(p)\)
yields the shown Bernoulli KL limit.
\end{proof}

\begin{corollary}[Two-point law of two-scale limit residual (bin-fold)]\label{cor:fold-bin-escort-two-scale-residual}
For any \(q\ge 0\) let \(W^{(q)}_m\sim\pi_{m,q}\), and define the scaled residual variable
\[
Z^{(q)}_m:=\Bigl(\frac{\varphi}{2}\Bigr)^m d^{\mathrm{bin}}_m\bigl(W^{(q)}_m\bigr).
\]
Then \(Z^{(q)}_m\Rightarrow Z^{(q)}\) and \(Z^{(q)}\) is a two-point distribution
\[
Z^{(q)}\in\{1,\varphi^{-1}\},
\qquad
\mathbb{P}\bigl(Z^{(q)}=\varphi^{-1}\bigr)=\theta(q),
\quad
\mathbb{P}\bigl(Z^{(q)}=1\bigr)=1-\theta(q).
\]
\end{corollary}

\begin{proof}
By Theorem \ref{thm:fold-bin-two-state-asymptotic},
\(
Z^{(q)}_m=\varphi^{-W^{(q)}_m(m)}+o(1)
\)
(in the sense of probability), and the limiting distribution of \(W^{(q)}_m(m)\) is given by Proposition \ref{prop:fold-bin-escort-lastbit}, so the conclusion follows.
\end{proof}

\begin{theorem}[Full Rényi divergence limit of bin-fold output distribution (against stable-layer uniform)]\label{thm:fold-bin-renyi-divergence-limit}
Let \(\mu_m(w):=d^{\mathrm{bin}}_m(w)/2^m\) be the pushforward output distribution of bin-fold, and \(u_m(w):=1/|X_m|\) the stable-layer uniform distribution.
For any \(\alpha>0\), define Rényi divergence \(D_\alpha(\mu_m\Vert u_m)\) (with \(\alpha=1\) taken as the continuous extension to KL divergence).
Then the limit
\[
D_\alpha^\infty:=\lim_{m\to\infty}D_\alpha(\mu_m\Vert u_m)
\]
exists and has a closed form:
\begin{enumerate}
  \item If \(\alpha\neq 1\), then
  \[
  D_\alpha^\infty=\frac{(2\alpha-2)\log\varphi+\log(\varphi+\varphi^{-\alpha})-\alpha\log\sqrt{5}}{\alpha-1}.
  \]
  \item If \(\alpha=1\), then
  \[
  D_1^\infty=\frac{\varphi}{\sqrt5}\log\Bigl(\frac{\varphi^2}{\sqrt5}\Bigr)+\frac{1}{\varphi\sqrt5}\log\Bigl(\frac{\varphi}{\sqrt5}\Bigr).
  \]
\end{enumerate}
Moreover, for each fixed \(\alpha\), the convergence rate is
\[
D_\alpha(\mu_m\Vert u_m)=D_\alpha^\infty+O\bigl((\varphi/2)^m\bigr).
\]
\end{theorem}

\begin{proof}
When \(\alpha\neq 1\), by definition
\[
D_\alpha(\mu_m\Vert u_m)
=\frac{1}{\alpha-1}\log\sum_{w\in X_m}\mu_m(w)^\alpha\,u_m(w)^{1-\alpha}.
\]
Substituting \(\mu_m(w)=d^{\mathrm{bin}}_m(w)/2^m\), \(u_m(w)=1/|X_m|\) yields
\[
D_\alpha(\mu_m\Vert u_m)
=\frac{1}{\alpha-1}\log\Bigl(|X_m|^{\alpha-1}\cdot 2^{-\alpha m}\cdot S_\alpha^{\mathrm{bin}}(m)\Bigr).
\]
By \(|X_m|=F_{m+2}\) and Proposition \ref{prop:fold-bin-power-sum-asymptotic},
\[
2^{-\alpha m}S_\alpha^{\mathrm{bin}}(m)
=\varphi^{-\alpha m}\Bigl(F_{m+1}+\varphi^{-\alpha}F_m\Bigr)\Bigl(1+O\bigl((\varphi/2)^m\bigr)\Bigr).
\]
Therefore
\[
D_\alpha(\mu_m\Vert u_m)
=\frac{1}{\alpha-1}\log\Bigl(F_{m+2}^{\alpha-1}\varphi^{-\alpha m}\bigl(F_{m+1}+\varphi^{-\alpha}F_m\bigr)\Bigr)
+O\bigl((\varphi/2)^m\bigr).
\]
Substituting Binet's asymptotics \(F_n=\varphi^n/\sqrt5+O(\varphi^{-n})\) and simplifying yields the limiting constant
\[
\frac{1}{\alpha-1}\Bigl((2\alpha-2)\log\varphi+\log(\varphi+\varphi^{-\alpha})-\alpha\log\sqrt5\Bigr),
\]
namely the shown \(D_\alpha^\infty\). The error bound follows from combining the relative error of Proposition \ref{prop:fold-bin-power-sum-asymptotic} with Binet's exponential remainder.
\par
When \(\alpha=1\), \(D_1(\mu_m\Vert u_m)=D_{\mathrm{KL}}(\mu_m\Vert u_m)\), whose limit is given by Proposition \ref{prop:fold-bin-kl-constant}; it can also be obtained directly from the two-state limit ratios \(\mu_m/u_m\Rightarrow\{\varphi^2/\sqrt5,\varphi/\sqrt5\}\), yielding the shown closed form. This completes the proof.
\end{proof}

\begin{proposition}[Universal collapse of \texorpdfstring{$f$}{f}-divergence (Csiszár index) for bin-fold]\label{prop:fold-bin-f-divergence-collapse}
Let \(f:(0,\infty)\to\RR\) be a convex function with \(f(1)=0\), and define the Csiszár \(f\)-divergence
\[
D_f(\mu_m\Vert u_m):=\sum_{w\in X_m}u_m(w)\,f\Bigl(\frac{\mu_m(w)}{u_m(w)}\Bigr).
\]
Then the limit \(D_f^\infty:=\lim_{m\to\infty}D_f(\mu_m\Vert u_m)\) exists and is given by a two-point closed form
\[
D_f^\infty=\frac{1}{\varphi}\,f\Bigl(\frac{\varphi^2}{\sqrt5}\Bigr)+\frac{1}{\varphi^2}\,f\Bigl(\frac{\varphi}{\sqrt5}\Bigr).
\]
Moreover, for fixed \(f\), the convergence rate is \(D_f(\mu_m\Vert u_m)=D_f^\infty+O\bigl((\varphi/2)^m\bigr)\).
\end{proposition}

\begin{proof}
Note that
\[
\frac{\mu_m(w)}{u_m(w)}=\frac{d^{\mathrm{bin}}_m(w)}{2^m}\cdot |X_m|
=d^{\mathrm{bin}}_m(w)\,\frac{F_{m+2}}{2^m}.
\]
By Theorem \ref{thm:fold-bin-two-state-asymptotic} and Binet's asymptotics, we obtain the uniform expansion
\[
\frac{\mu_m(w)}{u_m(w)}
=\frac{\varphi^{2-w_m}}{\sqrt5}+O\bigl((\varphi/2)^m\bigr),
\qquad (w\in X_m),
\]
and under the uniform measure \(u_m\),
\(
u_m(w_m=0)=F_{m+1}/F_{m+2}\to 1/\varphi
\),
\(
u_m(w_m=1)=F_m/F_{m+2}\to 1/\varphi^2
\).
Thus
\[
D_f(\mu_m\Vert u_m)
=\frac{F_{m+1}}{F_{m+2}}\,f\Bigl(\frac{\varphi^2}{\sqrt5}+O\bigl((\varphi/2)^m\bigr)\Bigr)
\;+\;\frac{F_m}{F_{m+2}}\,f\Bigl(\frac{\varphi}{\sqrt5}+O\bigl((\varphi/2)^m\bigr)\Bigr).
\]
By uniform continuity of convex functions on compact intervals, we may let \(m\to\infty\) to obtain the shown two-point limit; the error order similarly follows from the uniform expansion and \(F_{m+1}/F_{m+2}-1/\varphi=O(\varphi^{-2m})\). This completes the proof.
\end{proof}

\begin{proposition}[Identifiability of golden parameter: strict monotone inversion of power divergence limit constant]\label{prop:fold-bin-phi-identifiable-power-divergence}
For any \(\alpha>1\), take the strictly convex function
\[
f_\alpha(t):=t^\alpha-1-\alpha(t-1)\qquad(t>0),
\]
and denote its \(f\)-divergence limit by \(D_{f_\alpha}^\infty=\lim_{m\to\infty}D_{f_\alpha}(\mu_m\Vert u_m)\) (Proposition \ref{prop:fold-bin-f-divergence-collapse}).
Then the closed form is
\[
D_{f_\alpha}^\infty
=\frac{1}{5^{\alpha/2}}\Bigl(\varphi^{2\alpha-1}+\varphi^{\alpha-2}\Bigr)-1.
\]
Viewing the right-hand side as a function of parameter \(\xi>1\),
\[
\mathcal{D}_\alpha(\xi):=\frac{1}{5^{\alpha/2}}\Bigl(\xi^{2\alpha-1}+\xi^{\alpha-2}\Bigr)-1,
\]
then \(\mathcal{D}_\alpha:(1,\infty)\to\RR\) is strictly increasing; hence from the single scalar \(D_{f_\alpha}^\infty\) one can uniquely invert \(\xi\), thus in this system one can uniquely recover \(\xi=\varphi\).
\end{proposition}
\begin{proof}
By definition, \(\sum_{w\in X_m}u_m(w)\,\frac{\mu_m(w)}{u_m(w)}=\sum_w \mu_m(w)=1\), so for any \(m\),
\[
D_{f_\alpha}(\mu_m\Vert u_m)=\sum_{w}u_m(w)\Bigl(\frac{\mu_m(w)}{u_m(w)}\Bigr)^\alpha-1.
\]
Letting \(m\to\infty\), by the two-point limit law of Proposition \ref{prop:fold-bin-f-divergence-collapse}, we obtain
\[
D_{f_\alpha}^\infty
=\frac{1}{\varphi}\Bigl(\frac{\varphi^2}{\sqrt5}\Bigr)^\alpha+\frac{1}{\varphi^2}\Bigl(\frac{\varphi}{\sqrt5}\Bigr)^\alpha-1
=\frac{1}{5^{\alpha/2}}\Bigl(\varphi^{2\alpha-1}+\varphi^{\alpha-2}\Bigr)-1.
\]
For \(\xi>1\), the derivative of \(\mathcal{D}_\alpha\) is
\[
\mathcal{D}_\alpha'(\xi)=\frac{1}{5^{\alpha/2}}\Bigl((2\alpha-1)\xi^{2\alpha-2}+(\alpha-2)\xi^{\alpha-3}\Bigr).
\]
When \(\alpha\ge 2\), the above is clearly positive.
When \(1<\alpha<2\), since \(\xi^{\alpha+1}\ge 1\) and
\(
(2\alpha-1)-(2-\alpha)=3(\alpha-1)>0
\), we have
\[
(2\alpha-1)\xi^{2\alpha-2}+(\alpha-2)\xi^{\alpha-3}
\ge \xi^{\alpha-3}\bigl((2\alpha-1)\xi^{\alpha+1}-(2-\alpha)\bigr)
>0.
\]
Therefore \(\mathcal{D}_\alpha'(\xi)>0\) for all \(\xi>1\), so \(\mathcal{D}_\alpha\) is strictly increasing and hence invertible. This completes the proof.
\end{proof}

\begin{proposition}[Two-scale expansion of gauge volume density (bin-fold)]\label{prop:fold-bin-gauge-density-two-scale}
Let \(\mathcal{G}^{\mathrm{bin}}_m\) be the fold-gauge group volume of bin-fold
\(
|\mathcal{G}^{\mathrm{bin}}_m|=\prod_{w\in X_m} d^{\mathrm{bin}}_m(w)!
\)
and define the density \(g_m:=2^{-m}\log|\mathcal{G}^{\mathrm{bin}}_m|\).
Denote
\(
\kappa_m:=\mathbb{E}_{\mu_m}\bigl[\log d^{\mathrm{bin}}_m(W_m)\bigr]
\)
(the fiber information term in Proposition \ref{prop:fold-bin-kl-constant}).
Then as \(m\to\infty\), the two-scale expansion holds
\[
g_m=\kappa_m-1+\frac{(\varphi/2)^m}{2\sqrt5}\Bigl(\varphi^2 m\log\!\Bigl(\frac{2}{\varphi}\Bigr)+\varphi^2\log(2\pi)-\log\varphi\Bigr)+O\bigl((\varphi/2)^m\bigr).
\]
\end{proposition}

\begin{proof}
By Theorem \ref{thm:fold-bin-gauge-volume-stirling-second-order},
\[
g_m=\kappa_m-1+\frac{1}{2^{m+1}}\Bigl(|X_m|\log(2\pi)+\sum_{w\in X_m}\log d^{\mathrm{bin}}_m(w)\Bigr)+R_m,
\qquad 0\le R_m\le \frac{|X_m|}{12\cdot 2^m}.
\]
By the uniform logarithmic expansion of Theorem \ref{thm:fold-bin-two-state-asymptotic},
\[
\log d^{\mathrm{bin}}_m(w)=m\log\!\Bigl(\frac{2}{\varphi}\Bigr)-w_m\log\varphi+O\bigl((\varphi/2)^m\bigr),
\qquad (w\in X_m),
\]
summing and using
\(
\card\{w\in X_m:w_m=0\}=F_{m+1},\ \card\{w\in X_m:w_m=1\}=F_m
\)
yields
\[
\sum_{w\in X_m}\log d^{\mathrm{bin}}_m(w)
=m\log\!\Bigl(\frac{2}{\varphi}\Bigr)|X_m|-\log\varphi\cdot F_m+O\bigl(|X_m|(\varphi/2)^m\bigr).
\]
Substituting back and using Binet's asymptotics \(F_{m+2}=|X_m|=\varphi^{m+2}/\sqrt5+O(\varphi^{-m})\), \(F_m=\varphi^m/\sqrt5+O(\varphi^{-m})\) yields
\[
\frac{|X_m|}{2^{m+1}}
=\frac{\varphi^2}{2\sqrt5}\Bigl(\frac{\varphi}{2}\Bigr)^m+O\bigl((\varphi/2)^m\bigr),
\qquad
\frac{F_m}{2^{m+1}}
=\frac{1}{2\sqrt5}\Bigl(\frac{\varphi}{2}\Bigr)^m+O\bigl((\varphi/2)^m\bigr),
\]
and combining with \(R_m=O(|X_m|/2^m)=O((\varphi/2)^m)\) yields the shown expansion. This completes the proof.
\end{proof}

\begin{proposition}[Scale limit of bin-fold degeneracy spectrum: tail counts and two kinks in continuous capacity curve]\label{prop:fold-bin-degeneracy-tail-capacity-kinks}
For each \(m\ge 1\), define the bin-fold degeneracy tail count function
\[
N^{\mathrm{bin}}_m(t):=\bigl|\{w\in X_m:\ d^{\mathrm{bin}}_m(w)\ge t\}\bigr|,\qquad t\ge 0,
\]
and the continuous capacity curve
\[
\mathcal{C}^{\mathrm{bin,cont}}_m(T):=\sum_{w\in X_m}\min\bigl(d^{\mathrm{bin}}_m(w),T\bigr),\qquad T\ge 0.
\]
Then for all \(T\ge 0\), the strict identity holds
\[
\mathcal{C}^{\mathrm{bin,cont}}_m(T)=\int_{0}^{T}N^{\mathrm{bin}}_m(t)\,\dd t.
\]
For fixed \(s>0\), define the scaled processes
\[
\widehat N^{\mathrm{bin}}_m(s):=\frac{1}{|X_m|}\,N^{\mathrm{bin}}_m\!\Bigl(s\Bigl(\frac{2}{\varphi}\Bigr)^m\Bigr),
\qquad
\widehat C^{\mathrm{bin}}_m(s):=2^{-m}\,\mathcal{C}^{\mathrm{bin,cont}}_m\!\Bigl(s\Bigl(\frac{2}{\varphi}\Bigr)^m\Bigr).
\]
Then for each \(s\notin\{\varphi^{-1},1\}\), the pointwise limit holds
\[
\widehat N^{\mathrm{bin}}_m(s)\longrightarrow
\begin{cases}
1, & 0<s<\varphi^{-1},\\[2pt]
\varphi^{-1}, & \varphi^{-1}<s<1,\\[2pt]
0, & s>1,
\end{cases}
\]
and \(\widehat C^{\mathrm{bin}}_m\) converges uniformly on every compact set to the piecewise linear function
\[
\widehat C^{\mathrm{bin}}_\infty(s)
=\frac{\varphi}{\sqrt5}\min\{1,s\}+\frac{1}{\sqrt5}\min\{\varphi^{-1},s\}.
\]
Therefore in the bin-fold limit, the scaled "visible capacity-tail spectrum" exhibits two rigid kinks at \(s=\varphi^{-1}\) and \(s=1\).
\end{proposition}
\begin{proof}
The capacity-tail function identity follows from \(\min(d,T)=\int_0^T\ind_{\{t\le d\}}\,\dd t\), summing over \(w\) and interchanging summation and integration.
\par
Let \(t_m(s):=s(2/\varphi)^m\).
By the uniform asymptotics of Theorem \ref{thm:fold-bin-two-state-asymptotic},
\[
\Bigl(\frac{\varphi}{2}\Bigr)^m d^{\mathrm{bin}}_m(w)=\varphi^{-w_m}+O\bigl((\varphi/2)^m\bigr),
\qquad (w\in X_m),
\]
so for fixed \(s\notin\{\varphi^{-1},1\}\), on sufficiently large \(m\) the inequality
\(
d^{\mathrm{bin}}_m(w)\ge t_m(s)
\)
is equivalent to the last-bit condition:
\begin{itemize}
  \item If \(0<s<\varphi^{-1}\), then it holds for all \(w\in X_m\);
  \item If \(\varphi^{-1}<s<1\), then if and only if \(w_m=0\);
  \item If \(s>1\), then it holds for no \(w\in X_m\).
\end{itemize}
Therefore
\[
N^{\mathrm{bin}}_m(t_m(s))
=
\begin{cases}
|X_m|,& 0<s<\varphi^{-1},\\
|A_0|,& \varphi^{-1}<s<1,\\
0,& s>1,
\end{cases}
\qquad A_0=\{w:\ w_m=0\},
\]
thus
\(
\widehat N^{\mathrm{bin}}_m(s)\to 1
\)
or
\(
\widehat N^{\mathrm{bin}}_m(s)\to |A_0|/|X_m|=F_{m+1}/F_{m+2}\to \varphi^{-1}
\), yielding the shown pointwise limit.
\par
For the capacity curve, partitioning \(X_m\) by \(w_m\in\{0,1\}\) and using the two-state uniform asymptotics yields
\[
\mathcal{C}^{\mathrm{bin,cont}}_m(t_m(s))
=\sum_{w_m=0}\min\!\Bigl(\Bigl(\frac{2}{\varphi}\Bigr)^m+O(1),\ s\Bigl(\frac{2}{\varphi}\Bigr)^m\Bigr)
\;+\;
\sum_{w_m=1}\min\!\Bigl(\varphi^{-1}\Bigl(\frac{2}{\varphi}\Bigr)^m+O(1),\ s\Bigl(\frac{2}{\varphi}\Bigr)^m\Bigr).
\]
For fixed \(s\) and sufficiently large \(m\), the leading order of the two \(\min\) terms are respectively \((2/\varphi)^m\min\{1,s\}\) and \((2/\varphi)^m\min\{\varphi^{-1},s\}\), with errors of \(O(|X_m|)\).
Therefore
\[
\widehat C^{\mathrm{bin}}_m(s)
=2^{-m}\Bigl(\frac{2}{\varphi}\Bigr)^m\Bigl(|A_0|\min\{1,s\}+|A_1|\min\{\varphi^{-1},s\}\Bigr)+O\!\Bigl(\frac{|X_m|}{2^m}\Bigr),
\]
where \(|A_0|=F_{m+1}\), \(|A_1|=F_m\).
By Binet's asymptotics \(F_{m+1}/\varphi^m\to \varphi/\sqrt5\), \(F_m/\varphi^m\to 1/\sqrt5\), and noting
\(
2^{-m}(2/\varphi)^m=\varphi^{-m}
\), we obtain
\[
\widehat C^{\mathrm{bin}}_m(s)\to \frac{\varphi}{\sqrt5}\min\{1,s\}+\frac{1}{\sqrt5}\min\{\varphi^{-1},s\}.
\]
The error term \(O(|X_m|/2^m)=O((\varphi/2)^m)\) is uniform on every compact set, thus yielding uniform convergence. This completes the proof.
\end{proof}

\begin{corollary}[Endogenous recovery of \texorpdfstring{\(2\pi\)}{2pi} from fold-gauge volume and fiber information term]\label{cor:fold-bin-recover-2pi}
Adopting the notation of \(g_m,\kappa_m\) from Proposition \ref{prop:fold-bin-gauge-density-two-scale},
define the constant estimate sequence given directly by folding data
\[
C_m
:=\exp\!\Biggl(
\frac{2^{m+1}}{|X_m|}\bigl(g_m-\kappa_m+1\bigr)
-\frac{1}{|X_m|}\sum_{w\in X_m}\log d^{\mathrm{bin}}_m(w)
\Biggr).
\]
Then the limit exists and satisfies
\[
\lim_{m\to\infty}C_m=2\pi.
\]
Moreover, the convergence has exponential rate: there exists a constant \(c>0\) such that for all sufficiently large \(m\),
\[
\bigl|\log C_m-\log(2\pi)\bigr|\le c\Bigl(\frac{\varphi}{2}\Bigr)^m.
\]
\end{corollary}
\begin{proof}
For any integer \(d\ge 1\), Stirling's formula with remainder gives
\[
\log(d!)=d\log d-d+\frac12\log(2\pi d)+\varepsilon(d),
\qquad |\varepsilon(d)|\le \frac{1}{12d}.
\]
Summing over \(d=d^{\mathrm{bin}}_m(w)\) and dividing by \(2^m\), using \(\sum_w d^{\mathrm{bin}}_m(w)=2^m\), we obtain
\[
g_m
=\kappa_m-1+\frac{1}{2^{m+1}}\sum_{w\in X_m}\log(2\pi d^{\mathrm{bin}}_m(w))
\;+\;\frac{1}{2^m}\sum_{w\in X_m}\varepsilon\bigl(d^{\mathrm{bin}}_m(w)\bigr).
\]
Substituting
\(
\sum_w \log(2\pi d)=|X_m|\log(2\pi)+\sum_w \log d
\)
and rearranging yields
\[
\log C_m=\log(2\pi)+\frac{2}{|X_m|}\sum_{w\in X_m}\varepsilon\bigl(d^{\mathrm{bin}}_m(w)\bigr).
\]
By the uniform asymptotics of Theorem \ref{thm:fold-bin-two-state-asymptotic}, there exist \(m_0\) and constant \(c_0>0\) such that for \(m\ge m_0\) and all \(w\in X_m\),
\(
d^{\mathrm{bin}}_m(w)\ge c_0(2/\varphi)^m
\).
Therefore
\[
\Bigl|\frac{2}{|X_m|}\sum_{w}\varepsilon\bigl(d^{\mathrm{bin}}_m(w)\bigr)\Bigr|
\le \frac{2}{|X_m|}\sum_{w}\frac{1}{12\,d^{\mathrm{bin}}_m(w)}
\le \frac{1}{6}\cdot \frac{1}{\min_w d^{\mathrm{bin}}_m(w)}
\le c\Bigl(\frac{\varphi}{2}\Bigr)^m,
\]
thus \(\log C_m\to\log(2\pi)\) with exponential convergence bound. This completes the proof.
\end{proof}

\begin{theorem}[Stirling-Bernoulli hierarchy of folding gauge group: endogenous readout of all Bernoulli coefficients from folding data]\label{thm:fold-bin-gauge-constant-stirling-bernoulli-hierarchy}
Adopting the definition of \(C_m\) from Corollary \ref{cor:fold-bin-recover-2pi}, and denoting \(\log C_m\) as its natural logarithm,
there exists a complete asymptotic expansion: for any fixed integer \(R\ge 1\), as \(m\to\infty\),
\[
\boxed{\;
\log C_m
=\log(2\pi)+
\sum_{r=1}^{R}
\frac{B_{2r}}{r(2r-1)}
\Bigl(\varphi^{-1}+\varphi^{2r-3}\Bigr)
\Bigl(\frac{\varphi}{2}\Bigr)^{(2r-1)m}
\;+\;
O\!\Bigl(\Bigl(\frac{\varphi}{2}\Bigr)^{(2R+1)m}\Bigr)\; } ,
\]
where \(B_{2r}\) are Bernoulli numbers.
Therefore, for each \(r\ge 1\), the Bernoulli coefficient \(B_{2r}\) can be uniquely recovered in the limit from the folding data sequence \(\{\log C_m\}_{m\ge 1}\) via successive elimination;
and combining with the standard identity
\[
B_{2r}=(-1)^{r-1}\frac{2(2r)!}{(2\pi)^{2r}}\zeta(2r),
\]
all \(\zeta(2r)\) can be viewed as asymptotic spectral information of this folding gauge constant.
\end{theorem}
\begin{proof}
Adopt the Stirling-Bernoulli asymptotic expansion (as an external tool): for any fixed \(R\ge 1\) and integer \(d\ge 1\),
\[
\log(d!)=d\log d-d+\frac12\log(2\pi d)+
\sum_{r=1}^{R}\frac{B_{2r}}{2r(2r-1)}\,d^{-(2r-1)}
\;+\;
O\!\bigl(d^{-(2R+1)}\bigr).
\]
Applying this to \(d=d^{\mathrm{bin}}_m(w)\) and summing over \(w\in X_m\), then dividing by \(2^m\), and organizing according to
\(
g_m:=2^{-m}\sum_w\log(d^{\mathrm{bin}}_m(w)!)
\)
and
\(
\kappa_m:=2^{-m}\sum_w d^{\mathrm{bin}}_m(w)\log d^{\mathrm{bin}}_m(w)
\), we obtain
\[
g_m
=\kappa_m-1+\frac{1}{2^{m+1}}\sum_{w\in X_m}\log(2\pi d^{\mathrm{bin}}_m(w))
\;+\;
\sum_{r=1}^{R}\frac{B_{2r}}{2r(2r-1)}\cdot
\frac{1}{2^{m}}\sum_{w\in X_m} d^{\mathrm{bin}}_m(w)^{-(2r-1)}
\;+\;
O\!\left(\frac{1}{2^{m}}\sum_{w\in X_m} d^{\mathrm{bin}}_m(w)^{-(2R+1)}\right).
\]
Substituting
\(
\sum_w\log(2\pi d)=|X_m|\log(2\pi)+\sum_w\log d
\)
and substituting into the definition of \(\log C_m\) (Corollary \ref{cor:fold-bin-recover-2pi}), the leading terms strictly cancel, leaving only the uniform average of the Bernoulli series:
\[
\log C_m
=\log(2\pi)+
\sum_{r=1}^{R}\frac{B_{2r}}{r(2r-1)}\cdot
\frac{1}{|X_m|}\sum_{w\in X_m} d^{\mathrm{bin}}_m(w)^{-(2r-1)}
\;+\;
O\!\left(\frac{1}{|X_m|}\sum_{w\in X_m} d^{\mathrm{bin}}_m(w)^{-(2R+1)}\right).
\]
\par
By the uniform asymptotics of Theorem \ref{thm:fold-bin-two-state-asymptotic},
\(
d^{\mathrm{bin}}_m(w)=\varphi^{-w_m}(2/\varphi)^m+O(1)
\)
and last-bit partitioning \(\#\{w\in X_m:w_m=0\}=F_{m+1}\), \(\#\{w\in X_m:w_m=1\}=F_m\) with \(F_{m+1}/F_{m+2}\to\varphi^{-1}\), \(F_m/F_{m+2}\to\varphi^{-2}\), we obtain for each fixed \(r\ge 1\),
\[
\frac{1}{|X_m|}\sum_{w\in X_m} d^{\mathrm{bin}}_m(w)^{-(2r-1)}
=\Bigl(\varphi^{-1}+\varphi^{2r-3}\Bigr)\Bigl(\frac{\varphi}{2}\Bigr)^{(2r-1)m}
\;+\;
O\!\Bigl(\Bigl(\frac{\varphi}{2}\Bigr)^{(2r+1)m}\Bigr).
\]
Similarly,
\(
|X_m|^{-1}\sum_w d^{-(2R+1)}
=O((\varphi/2)^{(2R+1)m})
\).
Substituting back yields the shown expansion and remainder order. The final two sentences about "successive elimination recovering \(B_{2r}\)" and the encoding of \(\zeta(2r)\) are standard corollaries: since the exponential orders \((\varphi/2)^{(2r-1)m}\) are mutually separated, one can uniquely read out each coefficient by successively removing lower-order terms and letting \(m\to\infty\). This completes the proof.
\end{proof}

\endinput


\begin{remark}[Two-state compression as analytic benchmark solution]\label{rem:fold-bin-solvable-model}
In the bin-fold case, Propositions \ref{prop:fold-bin-power-sum-asymptotic}--\ref{prop:fold-bin-gauge-density-two-scale} show that: all-order collision moments, full-temperature escort families, Rényi/\(f\)-divergence families, and the first visible correction term of gauge volume density can all be closed-form expressed under the same pair of two-state constants with uniform error order.
This structure provides an analytic benchmark for comparison with the more general folding mechanisms in the main text.
\end{remark}
